\section{本章小结}
\label{sec:graph-objects-summary}

本章围绕三个问题建立了图的基本语言：连接是什么，沿连接能到哪里，结构怎样记录和比较。图$\mathcal G=(V,E)$只记录有限个对象之间的直接连接，与纸上的画法无关。方向、边权和属性各自附加一类信息：方向决定哪些顶点序列可以通行，边权的含义由问题决定，长度相加、容量限制通过量、相似度鼓励两端接近；零权边仍是边，“能否到达”查边集，“累计多少”查边权。子图与诱导子图区分了“任意保留一部分边”和“保留选定顶点之间的全部边”。

沿连接能到哪里，由行走和路径回答。行走允许重复，便于拼接与计数；路径不重复顶点，对应实际路线；任何行走中都含有一条同端点的路径，所以判断可达时两者等价。无向图中可达关系是等价关系，把顶点集唯一地划分为连通分量；有向图中可达不对称，需要区分要求双向到达的强连通与忽略方向的弱连通，前者蕴含后者而反之不然。分隔集、割点和桥刻画了删除顶点或边对连接的影响，判断的依据始终是有没有替代路径：一条边是桥，当且仅当它不在任何环上。树是连通而无环的图，等价于任意两点之间路径唯一；$n$个顶点的树恰有$n-1$条边，连通且有$n-1$条边的图是树，而$n$个顶点、$c$个分量的生成森林恰有$n-c$条边。

同一张图可以用边列表、邻接表、邻接矩阵和关联矩阵记录，四者保存的连接相同，支持的访问和代价不同；邻接矩阵还把图的组合问题与矩阵运算联系起来，$\bm A^k$的条目计数长度为$k$的行走，但其中包含重复顶点。同构用保持邻接与不相邻的双射刻画“只差一套名称”，在矩阵上表现为置换共轭$\bm A'=\bm P\bm A\bm P^{\top}$；度、分量数等不变量可以排除同构，却不足以确认同构。拿到一张图时，可以用这些表示和不变量核对：所关心的结构是否被保留，附加的数值和方向是否交代清楚。

本章定义了连通分量、强连通分量、生成树、有向无环图上的依赖和同构，却还没有给出高效求出它们的方法。第~\ref{chap:graph-traversal}章在邻接表上用广度优先和深度优先遍历实际计算这些对象，用拓扑排序安排有向无环图的次序，并用顶点标记细化比较图的结构；那里的算法输出，仍要按本章的定义与条件检验。
